% Final Q3 scientific wording; standalone fragment for review.
% Not compiled and not automatically inserted into the historical template.
\section{Behavioural responses to sound playback}
\paragraph{Operational response definitions.}
We retained all 278 records from 139 paired experiments and used the original
head-up, vocal and approach flags to define the eight joint response states.
In this dataset, the vocal flag is exactly equivalent to
$N_{\mathrm{after}}>N_{\mathrm{before}}$.
It therefore differs from the presence of any recorded call. The zero state
means that none of the original flags is positive; it does not establish
absence of calling or perception. The original annotation procedure remains
unconfirmed. Earlier analyses excluding records with a zero flag and a recorded
latency are retained as population-selection sensitivity analyses.

\paragraph{Response pattern and intensity.}
Within the supported type-specific matched reference populations, Flat, Mod and
TV playback were associated with increases of 40.25, 39.82 and 47.02 percentage
points in the probability of any original response flag, respectively.
Their pointwise 95\% experiment-cluster bootstrap intervals were
[28.84,52.34], [27.08,53.29] and [36.26,57.71] percentage points.
The corresponding adjusted differences in expected group call counts were
1.178 [0.764,1.652], 1.706 [0.830,2.706] and 1.481 [1.060,2.000] calls per
playback. These comparisons use different type-specific reference populations
of 30, 20 and 39 pairs, so their magnitudes do not directly establish a ranking
of sound types. Count predictions remain somewhat low in the upper tail.
The original fixed NB2 penalty was retained after the limited sensitivity test.
True behavioural duration cannot currently be identified from the unresolved
60-second encoding.

\paragraph{Window probabilities and recorded response times.}
Latency window probabilities are the marginals of the same joint flag model.
Head-up and vocal-label probabilities were higher for all three natural sounds.
Approach differences were positive for Mod and TV. The Flat approach contrast
was 5.33 [0.72,12.10] percentage points in the original-flag analysis, but its
interval crossed zero in the earlier filtered-population sensitivity analysis;
this result is therefore sensitive to the analysis rule.
Conditional recorded times were modelled separately using call type only.
There was no reliable evidence that natural playback generally shortened these
times. The Control condition had no positive Approach time observations;
conditional time predictions there borrow model information and are not
observed Control response times.

For vocal behaviour, the reported cumulative quantity is
\[
 F_V(t\mid x)=P\{E_V=1,\ L_V^{\mathrm{recorded}}\leq t\mid x\},
 \qquad E_V=I(N_{\mathrm{after}}>N_{\mathrm{before}}).
\]
It describes a final positive vocal-increase label together with a recorded
vocal time no greater than $t$. It is not the distribution of the first time
the cumulative number of calls exceeds baseline. The final increase criterion
uses the entire observation window. Likewise, recorded zero latencies are not
assumed to be exact instantaneous biological responses.

\paragraph{Cross-module compatibility and downstream use.}
Let $q=P(N_{\mathrm{after}}>0\mid x)$ and $p_V=P(E_V=1\mid x)$.
The event definitions imply $p_V\leq q$, with equality when baseline count is
zero. Separately fitted models did not enforce these constraints. All 36
representative scenarios used a zero baseline; ten had $p_V>q$, with a maximum
excess of 6.93 percentage points. The mean absolute discrepancy was 1.75
percentage points. The NB2 count distribution also supplies an independent
estimate of the same increase event, using its truncated upper tail.
No probabilities were calibrated or replaced. The response and count modules
may be reported separately, but they do not constitute a coherent joint
generative model for behaviour flags and call counts. Q4 must not multiply
these probabilities or infer new joint events without addressing this limitation.

All intervals are exploratory, pointwise, and conditional on fixed model and
analysis rules. Grouped prediction validation does not establish causality or
resolve unknown annotation mechanisms. The actual within-pair interval,
underlying recording identities and duration coding remain unresolved.
